\section*{Lowpass VCVS}
\noindent\colorbox{orange}{
  \begin{minipage}{1.0\linewidth}
    \paragraph{Richiesta (a)}
    Dimostrare, sulla base di semplici argomenti di natura qualitativa (ovvero prescindendo
    dall’espressione analitica della funzione di trasferimento) che il circuito proposto svolga
    effettivamente la funzione di filtro passa-basso, discutendone la risposta ai limiti del dominio
    delle frequenze reali e verificare il limite del guadagno a bassa frequenza ottenuto al punto
    1c.
  \end{minipage}
}

\paragraph{Risposta (a)}
\begin{itemize}
\item
  Per $f \to 0$ i condensatori $C_1$ e $C_2$ aprono i loro rami. Il circuito opera
  come un amplificatore \underline{non}-invertente. La tensione al nodo A \'e la
  tensione sorgente \textit{partita} dal partitore di tensione $(R_1, R_2)$
  \begin{equation*}
    V_A = \frac12 V_s
  \end{equation*}
  $R_3$ \'e una resistenza finita trascurabile rispetto all'impendenza ``grande'' dell'opamp,
  di conseguenza la tensione dell'ingresso invertente
  \begin{equation*}
    V_+ = \frac12 V_s
  \end{equation*}
  In maniera del tutto analoga a quanto detto in precedenza, la tensione dell'ingresso
  invertente
  \begin{equation*}
    V_- = \frac12 V_{OUT}
  \end{equation*}
  Se l'opamp opera in regime di linearit\'a $V_{OUT} = A_dV_d$
  \begin{equation*}
    V_{OUT} = \frac12 A_d(V_s - V_{OUT}) \implies \colorbox{blu}{\boxed{V_{OUT} = \frac{A_d}{2 + A_d}V_s}}
  \end{equation*}
  nel limite di basse frequenze il circuito ha un andamento compatibile con quello di
  un circuito filtro passa basso.
\item Per $f \to +\infty$ il condensatore $C_1$ ``corto-circuita'' il proprio ramo
  cos\'i che $V_+ = 0$. Se l'opamp opera in regime di linearit\'a $V_{OUT} = A_dV_d$
  \begin{equation*}
    V_{OUT} = A_d(0 - V_-) = -A_dV_{OUT} \implies \colorbox{blu}{\boxed{V_{OUT} = \frac{0}{1 + A_d} = 0}}
  \end{equation*}
  nel limite di alte frequenze il circuito ha un andamento compatibile con quello di
  un circuito filtro passa basso.
\end{itemize}

\noindent\colorbox{orange}{
  \begin{minipage}{1.0\linewidth}
    \paragraph{Richiesta (b)}
    Derivare analiticamente la trasformata di Laplace della funzione di trasferimento $A_v(s) =
    V_{OUT}(s)/V_S(s)$ nell’ipotesi che l’operazionale sia ideale ed i valori di resistenze e capacit\'a
    siano pari a quelli nominali, ovvero valgano le relazioni C1=C2=C, R1=R2=R, R3=2R,
    R4=R5 (\textit{suggerimento: si determini la tensione nel punto A del circuito in termini di VS e
      $V_{OUT}$; si determini quindi la tensione all’ingresso non invertente dell’operazionale e
      l’ulteriore relazione tra $V_A$ e $V_{OUT}$ assumendo che l'operazionale funzioni in regime lineare}).
  \end{minipage}
}

\paragraph{Risposta (b)}
La richiesta porta con se un suggerimento, seguendolo per il nodo A vale
\begin{equation*}
  \frac{V_S - V_A}{R} = \frac{V_A}{R} + \frac{V_A - V_B}{2R} + sC(V_A - V_{OUT})
\end{equation*}
quindi l'espressione della KCL per il nodo A include i termini $V_A$ e $V_{OUT}$
successivamente, sempre seguendo le istruzioni date si trovano ulteriori
relazioni per i nodi $V_A$ e $V_B$
\begin{equation*}
  \frac{V_A - V_B}{2R} = sCV_B
\end{equation*}
assumendo che l'opamp sia ideale e che operi in regime lineare si applica
il principio di corto-circuito virtuale $V_+ = V_-$ da cui
\begin{equation*}
  \begin{gathered}
    \begin{cases}
      \frac{V_{OUT} - V_-}{R} = \frac{V_-}{R} \\
      V_B = V_+ = V_-
    \end{cases} \implies
    \frac{V_{OUT} - V_B}{R} = \frac{V_B}{R} \\
    V_B = \frac12V_{OUT}
  \end{gathered}
\end{equation*}
ora $V_B$ \'e puramente espressione di $V_{OUT}$ questo permetter\'a
di trovare una relazione analoga per $V_A$
\begin{equation*}
  \begin{gathered}
    \begin{cases}
      \frac{V_A - V_B}{2R} = sCV_B \\
      V_B = \frac12V_{OUT}
    \end{cases} \implies
    \frac{V_A - \frac12V_{OUT}}{2R} = sC\frac12V_{OUT} \\
    V_A = \frac12(1 + 2sRC)V_{OUT}
  \end{gathered}
\end{equation*}
ora $V_B$ e $V_A$ sono espressioni unicamente di $V_{OUT}$. Ora si riduce
la prima equazione nodale trovata all'inizio di questa risposta
\begin{equation*}
  V_S = \frac12(5 + 2sRC)V_A - \frac12V_B - sRCV_{OUT}
\end{equation*}
e si sostituiscono $V_A$ e $V_B$ con le loro espressioni in termini di $V_{OUT}$
\begin{equation*}
  V_S = \frac14\left[(5 + 2sRC)(1 + 2sRC) - 1 - 4sRC\right]V_{OUT}
\end{equation*}
l'espressione finale per la funzione di trasferimento complessa
\begin{equation}
  \label{eq:lp-vcvs-tfunc}
  \colorbox{blu}{\boxed{H(s) = \frac{1}{(1 + RCs)^2}}}
\end{equation}

\noindent\colorbox{orange}{
  \begin{minipage}{1.0\linewidth}
    \paragraph{Richiesta (c)}
    Dalla funzione di trasferimento dedurre le caratteristiche del filtro osservate al punto 1)
    (guadagno a centro-banda, frequenza naturale, andamento di modulo/fase in funzione della
    frequenza) e confrontarli con quanto misurato ai punti 1c, 1d, 1e.
  \end{minipage}
}
\paragraph{Risposta (c)}
Le parti reali e immaginarie del numeratore e denominatore
\begin{equation*}
  \begin{cases}
    \mathcal{R}_{\omega}^N = 1 \\
    \mathcal{I}_{\omega}^N = 0 \\
    \mathcal{R}_{\omega}^D = 1 - R^2C^2\omega^2 \\
    \mathcal{I}_{\omega}^D = 2RC\omega
  \end{cases} \implies
  G(\omega) = \sqrt{\frac{1}{(1 - R^2C^2\omega^2)^2 + (2RC\omega)^2}}
\end{equation*}
il guadagno \'e ottenuto dall'espressione \ref{eq:da-tfunc-a-gain}
\begin{equation}
  \label{eq:lp-vcvs-gain}
  G(f) = \frac{1}{1 + 4\pi^2R^2C^2f^2}
\end{equation}
il denominatore \'e positivo monotono crescente, limitato inferiormente a $G(f) = 1$,
di conseguenza il guadagno massimo, essando il reciproco del denominatore \'e proprio
quell'estremo inferiore
\begin{equation}
  \label{eq:lp-vcvs-gain-max}
  \colorbox{blu}{\boxed{\max_{f \in [0, +\infty)} G(f) = G(0) = 1}}
\end{equation}
la frequenza centrale \'e la frequenza per cui il guadagno \'e massimo
\begin{equation}
  \label{eq:lp-vcvs-frequenza-centrale}
  \colorbox{blu}{\boxed{f = 0}}
\end{equation}
ad ogni polo della funzione di trasferimento complessa $H(s)$ corrisponde una
frequenza naturale\footnotemark

\footnotetext{
  Il viceversa non \'e vero in generale, alcune frequenze naturali
  corrispondono a particolari condizioni iniziali
  ``Desoer, Charles (1969). Basic circuit theory. McGraw-Hill pag. 643''
}

\begin{equation}
  \label{eq:lp-vcvs-frequenza-naturale}
  \colorbox{blu}{\boxed{f = \frac{1}{2\pi RC}}}
\end{equation}

Per calcolare l'espressione della fase bisogna prima identificare il semi-piano di appartenenza,
il sengo della parte reale \'e determinata da
\begin{equation*}
  1-R^2C^2\omega^2
\end{equation*}
il segno non \'e be determinato per ogni $\omega \in [0, +\infty)$. Invece, il segno della parte
immaginaria \'e determinata da
\begin{equation*}
  2RC\omega > 0
\end{equation*}
di conseguenza il semipiano \'e quello \textit{immaginario positivo} da cui la fase $\phi$ in
funzione della frequenza viene calcolata attraverso la \ref{eq:da-tfunc-a-fase-immaginario-positivo}
e la sostituzione $\omega \to 2\pi f$
\begin{equation}
  \label{eq:lp-vcvs-fase}
  \colorbox{blu}{\boxed{\phi(f) = -\frac{\pi}{2} + \arctan\left(\frac{1 - 4\pi^2R^2C^2f^2}{4\pi RCf}\right)}}
\end{equation}

\noindent\colorbox{orange}{
  \begin{minipage}{1.0\linewidth}
    \paragraph{Richiesta (d)}
    Dare un’interpretazione qualitativa e quantitativa delle caratteristiche del segnale di uscita
    osservato al punto 1f. Che tipo di segnale appare? A quale ampiezza?
  \end{minipage}
}
